\section*{第 \ref{ch:b3:remote-sensing-gis} 章——遥感与地理信息系统}

\begin{solution}{exo:b3:remote-sensing-gis:1}
农田：$(0.42 - 0.06)/(0.42 + 0.06) = 0.75$，茂密的绿色植被，作物
处于满叶期。耕地：$(0.28 - 0.21)/(0.28 + 0.21) = 0.14$，几乎
裸露的地面，土壤在红外方向上略有增亮。
\end{solution}

\begin{solution}{exo:b3:remote-sensing-gis:2}
(a)~$185000/30 \approx 6170$ 个像元；$185000/15 \approx 12300$ 个
全色波段中的像元。(b)~$2^{12} = 4096$ 个灰度级；$2^8 = 256$ 个。(c)~单颗卫星时 16 天
；若 Landsat~9 与之同轨、相隔 8 天，则为 8 天。
\end{solution}

\begin{solution}{exo:b3:remote-sensing-gis:3}
(a)~一个条纹为 $\lambda/2 = \qty{27.7}{mm}$；六个条纹
对应\omterm{def:b3:remote-sensing-gis:insar}{视线方向}上 \qty{166}{mm} 的变化，此处远离卫星，因为
盆地在下沉。(b)~$56.6/2 = \qty{28.3}{mm}$。
\end{solution}

\begin{solution}{exo:b3:remote-sensing-gis:4}
(a)~$t = 2d/c = 2\times\num{5e5}/\num{3.00e8} = \qty{3.34}{ms}$.
(b)~$c\,\delta t/2 = \num{3.00e8}\times\num{8e-10}/2 = \qty{0.12}{m}$；由许多光子采样的表面
定位精度高得多，精度
与光子数的平方根成反比。
\end{solution}

\begin{solution}{exo:b3:remote-sensing-gis:5}
(a)~按国际单位制，$L = \qty{9.0e6}{W.m^{-2}.sr^{-1}.m^{-1}}$。取 $c_1 =
2hc^2 = \qty{1.191e-16}{W.m^2.sr^{-1}}$，$c_2 = \qty{0.014388}{m.K}$，
$\lambda^5 = (\num{1.09e-5})^5 = \num{1.539e-25}$, $c_1/(\lambda^5L) = 86.0$,
则 $T_b = c_2/(\lambda\ln 87.0) = \qty{295.6}{K}$。(b)~动力学温度
满足 $0.98\,B_\lambda(T) = L$，即 $B_\lambda(T) = L/0.98$：同一
公式给出 $T = \qty{296.9}{K}$，高出 \qty{1.3}{K}。
\end{solution}

\begin{solution}{exo:b3:remote-sensing-gis:6}
(a)~$R = 693/\cos\ang{35} = \qty{846}{km}$。(b)~真实天线：$R\lambda/L =
\num{8.46e5}\times0.0555/12.3 = \qty{3.8}{km}$；合成孔径：$L/2 =
\qty{6.2}{m}$，分辨率约高 600 倍。(c)~合成孔径为
某点停留在波束内的航迹长度，$R\lambda/L =
\qty{3.8}{km}$.
\end{solution}

\begin{solution}{exo:b3:remote-sensing-gis:7}
$\partial z/\partial x = (418 - 405)/60 = 0.217$（东减西），
$\partial z/\partial y = (418 - 401)/60 = 0.283$（北减南）；
$|\nabla z| = 0.357$，坡度为 $\arctan 0.357 = \ang{19.6}$。地面
向东北升高，因此坡面朝西南：$A =
\operatorname{atan2}(-0.217, -0.283) = \ang{217}$。山体阴影：$\cos\ang{45}
\cos\ang{19.6} + \sin\ang{45}\sin\ang{19.6}\cos(\ang{315} - \ang{217}) =
0.666 - 0.033 = 0.63$，比平坦地面（0.71）略暗：坡面
略微背离西北方向的光源。
\end{solution}

\begin{solution}{exo:b3:remote-sensing-gis:8}
行合计 52、48、38、22；列合计 54、46、38、22；对角线之和为 133，
样本总数 160，总体精度 \qty{83.1}{\percent}。用户精度：水体
$48/52 = 0.92$，森林 $37/48 = 0.77$，作物 $29/38 = 0.76$，城镇
$19/22 = 0.86$。生产者精度：水体 $48/54 = 0.89$，森林 $37/46 = 0.80$，
作物 $29/38 = 0.76$，城镇 $19/22 = 0.86$。$p_e = (52\cdot54 + 48\cdot46 +
38\cdot38 + 22\cdot22)/160^2 = 0.271$, $\kappa = (0.831 - 0.271)/(1 -
0.271) = 0.77$。作物最常被误分：38 个真实作物像元中有 9 个
被分到了其他类别，多数被分为森林。
\end{solution}

\begin{solution}{exo:b3:remote-sensing-gis:9}
(a)~$L = 0.04\,B_\lambda(1443) + 0.96\,B_\lambda(290)$，在
\qty{10.9}{\micro\metre} 处为 \qty{28.6}{W.m^{-2}.sr^{-1}.\micro\metre^{-1}}，
其\omterm{def:b3:remote-sensing-gis:brightness}{亮度温度}为 \qty{396}{K}，而平均温度为
$0.04\times1443 + 0.96\times290 = \qty{336}{K}$。(b)~在
\qty{2.2}{\micro\metre} 处，$B_\lambda(1443)/B_\lambda(290) = \num{6.8e7}$：
熔岩完全主导了该像元，其\omterm{def:b3:remote-sensing-gis:brightness}{亮度温度}为
\qty{845}{K}。短波红外波段对小热点的探测能力强得多；
热红外波段对整个像元敏感，它
测量的是其周围的地面。
\end{solution}

\begin{solution}{exo:b3:remote-sensing-gis:10}
卫星位于东侧时，$d_1 = u_E\sin\theta + u_z\cos\theta$；位于
西侧时，$d_2 = -u_E\sin\theta + u_z\cos\theta$。相加与相减，
$u_z = (d_1 + d_2)/(2\cos\theta) = 102/1.638 = \qty{62.3}{mm}$，并得 $u_E =
(d_1 - d_2)/(2\sin\theta) = 26/1.147 = \qty{22.7}{mm}$。单次过境只给出一个
关于两个未知量的方程：隆升与向东运动的任意组合，只要
在\omterm{def:b3:remote-sensing-gis:insar}{视线方向}上的投影相同，便给出相同的条纹。
\end{solution}

\begin{solution}{exo:b3:remote-sensing-gis:11}
(a)~由\cref{thm:b3:remote-sensing-gis:fringe}可知，$\Delta\phi =
4\pi\Delta r/\lambda < \pi$，要求 $\Delta r < \lambda/4 =
\qty{13.9}{mm}$。(b)~$u_z = k z_0(r^2 + z_0^2)^{-3/2}$，其中 $k = (1 -
\nu)\Delta V/\pi$; $\mathrm du_z/\mathrm dr = -3kz_0r(r^2 + z_0^2)^{-5/2}$,
其导数在 $(r^2 + z_0^2) = 5r^2$，即 $r = z_0/2$ 处为零。于是
\[
\left|\frac{\mathrm du_z}{\mathrm dr}\right| = \frac{3k}{2z_0^3}\Big(\frac54\Big)^{-5/2}
= \frac{0.859\,k}{z_0^3} = \frac{0.859\times\num{2.39e6}}{\num{6.4e10}} = \num{3.2e-5},
\]
\qty{32}{mm} 每千米。(c)~$13.9/32 = \qty{0.43}{km}$：像元应小于
约 \qty{430}{m}，Sentinel-1 的像元很容易满足这一要求。
\end{solution}

\begin{solution}{exo:b3:remote-sensing-gis:12}
(a)~设 $N, R \ge 0$ 不同时为零，则 $|N - R| \le N + R$，故 $|\mathrm{NDVI}|
\le 1$；将 $N$ 和 $R$ 换成 $aN$ 和 $aR$ 时，分子和
分母都乘以 $a$。(b)~未校正时：$(0.502 - 0.051)/(0.502 + 0.051) =
0.816$；校正后：$(0.512 - 0.081)/(0.512 + 0.081) = 0.727$。(c)~程辐射
使绿色植被的指数偏低，此处约低十分之一，且
偏低幅度随霾而变化：未经校正的影像所做的地图会
与植被一样反映天气，不同日期之间的比较
将毫无意义。
\end{solution}

\begin{solution}{pb:b3:remote-sensing-gis:1}
\textbf{1.} 森林 $(0.38 - 0.04)/0.42 = 0.81$；1890 年熔岩流 $0.23/0.37 =
0.62$；1960 年熔岩流 $0.10/0.32 = 0.31$；2005 年熔岩流 $0.02/0.26 = 0.08$。

\textbf{2.} 依次为 1890、1960、2005 年：最老的熔岩流最绿，
即绿度的排序与年龄的排序一致；它尚未到达森林。

\textbf{3.} 裸岩在近红外的反射率略高于红光，
与本章图中的棕色沙地相似；地衣和最早的
植物只略增一点：0.08 接近几乎裸露地面的值。

\textbf{4.} 背光的坡面使两个波段的亮度按相同比例变暗；
比值可以抵消这一影响，而仅用近红外反射率
则会把阴影中的森林误判为稀疏植被。

\textbf{5.} 由\cref{prop:b3:remote-sensing-gis:planck-inversion}，取 
$\lambda = \qty{10.9}{\micro\metre}$：高原处 \qty{291.1}{K}，
喷气孔像元 \qty{302.3}{K}。

\textbf{6.} 由 $0.97\,B_\lambda(T) = \num{8.40e6}$ 得 $T = \qty{293.0}{K}$。

\textbf{7.} $B_\lambda(\qty{370}{K}) =
\qty{22.48}{W.m^{-2}.sr^{-1}.\micro\metre^{-1}}$; $f = (9.95 - 8.40)/(22.48
- 8.40) = 0.110$，约 \qty{1100}{m^2}，占 \qty{e4}{m^2} 像元。

\textbf{8.} 辐射亮度是温度的凸函数
（见\cref{prop:b3:remote-sensing-gis:mixed}）：热部分所增加的辐射亮度
超过其在温度中所占的份额，因此像元看起来比其平均温度更高，
$0.110\times370 + 0.890\times291.1 = \qty{299.8}{K}$，而观测值为
\qty{302.3}{K}。

\textbf{9.} 每次回波的相位取决于精确路径；视点的较大变化
会改变散射体回波之和，使
\omterm{def:b3:remote-sensing-gis:insar}{干涉图}失去相干性，并引入来自地形的相位。
$48/12 = 4$ 个循环（单颗卫星）。

\textbf{10.} $5.4\times27.73 = \qty{149.7}{mm}$，朝向卫星。

\textbf{11.} $u_z = 149.7/\cos\ang{35} = \qty{182.8}{mm}$。

\textbf{12.} 地面既向上也向外移动；东侧的
向外运动朝向卫星，与隆升叠加；西侧
则相减（见\cref{prop:b3:remote-sensing-gis:mogi-los}）。形变
在西侧下降得更快，那里条纹较密；图样
向东伸展。

\textbf{13.} 沿该线，水平运动为南–北向，
与\omterm{def:b3:remote-sensing-gis:insar}{视线方向}垂直，因此雷达只能探测到 $u_z\cos\theta$，其形状与
隆升相同。隆升在 $r = 0.766\,z_0$ 处减半，因此 $z_0 =
3.4/0.766 = \qty{4.44}{km}$。

\textbf{14.} 5.4 的一半：2.7 个条纹。

\textbf{15.} 经时间间隔换算后两幅\omterm{def:b3:remote-sensing-gis:insar}{干涉图}上的形变相同，或形变
逐对持续增长，属真实形变；大气图样
随日期而变，且常随地形起伏。
另一图对若显示相同的靶心图样，则
证实该形变。

\textbf{16.} $\Delta V = \pi u_z(0)z_0^2/(1 - \nu) = \pi\times0.1828\times
4436^2/0.75 = \qty{1.51e7}{m^3}$。

\textbf{17.} 每天 $\num{1.51e7}/48 = \qty{3.1e5}{m^3}$，每年 \qty{1.1e8}{m^3}
若速率保持不变。

\textbf{18.} $u_r \propto r/(r^2 + z_0^2)^{3/2}$ 在
其导数为零处最大，即 $r^2 + z_0^2 = 3r^2$，$r = z_0/\sqrt2 = \qty{3.14}{km}$；
此处 $u_r = \qty{70}{mm}$，为中央隆升的 \qty{38}{\percent}。

\textbf{19.} 半径 \qty{1.48}{km}，体积 $\frac43\pi(1.48)^3 =
\qty{13.6}{km^3}$；变化量占其体积的 $\num{1.51e7}/\num{1.36e10} =
\qty{0.11}{\percent}$。

\textbf{20.} 不能：一个条纹的十分之一约占计数的 \qty{2}{\percent}，而
深度由半宽 \qty{3.4}{km} 读取，或许精确到
\qty{0.1}{km}，以平方形式计入，占 \qty{6}{\percent}；\qty{3}{\percent}
的修正量在该不确定度之内。

\textbf{21.} 喷气孔区随时间的升温（第二部分），
隆起下方小地震的增加、释放气体的变化，
以及用卫星定位或水准测量在
山坡上的若干点测得的隆升。

\textbf{22.} 哈尔弗伦山下的储库增长了
\textbf{1500 万~\unit[text-series-to-math=true]{m^3}}（$\num{1.51e7}\,\unit{m^3}$）。
\end{solution}
